%==============================================================================
\section{Introduction}
\label{sec:introduction}

Two-dimensional echocardiography is the most widely used cardiac imaging
modality worldwide: it is portable, inexpensive, non-ionising, and provides
real-time visualisation of cardiac structure and function. Its core
quantitative output --- left ventricular ejection fraction (LVEF), computed
from end-diastolic and end-systolic volumes via the biplane Simpson's method
\citep{lang2015} --- drives a large fraction of clinical cardiac decisions,
from diagnosis of heart failure to monitoring of cardiotoxic chemotherapy.
The clinical importance of LVEF is matched by the operator-dependent
nature of its measurement: published inter-observer LVEF variability is
approximately $7$--$10\%$ \citep{leclerc2019}, which directly limits the
sensitivity of LVEF as a clinical decision boundary.

Automated segmentation of the left ventricular endocardium (LV-endo),
epicardium (LV-epi), and left atrium (LA) from two-dimensional apical
2-chamber (2CH) and 4-chamber (4CH) views would, in principle, eliminate
this operator-dependent variability and provide reproducible volumetric
measurements. The technical challenge is substantial: ultrasound suffers
from low signal-to-noise ratio, speckle noise, low boundary contrast,
acoustic shadowing, and view-dependent geometric foreshortening.

\paragraph{Architectural pluralism.}
Three paradigms have emerged as competitive solutions in this space:
\begin{itemize}[leftmargin=*,nosep]
\item \textbf{Pure convolutional networks.} U-Net
      \citep{ronneberger2015} and its descendants (residual encoders, SE
      blocks, attention gates, deep supervision)
      \citep{milletari2016,oktay2018} remain the dominant architectural
      family, with nnU-Net \citep{isensee2021} serving as a strong
      auto-configured baseline.
\item \textbf{Hybrid CNN--Transformer networks.} TransUNet
      \citep{chen2021transunet} and similar hybrids combine a
      convolutional encoder with a Vision Transformer bottleneck,
      explicitly trading off local detail (CNN) against global context
      (self-attention).
\item \textbf{Pure Transformer networks.} Swin-UNet \citep{cao2022}
      replaces the convolutional \emph{backbone} with a hierarchy of
      windowed self-attention blocks --- retaining only a shallow
      convolutional patch-embedding stem --- demonstrating that
      predominantly attention-based encoders can match or exceed CNN
      performance on a range of medical tasks.
\end{itemize}

\paragraph{The benchmarking gap.}
Despite this architectural pluralism, the literature offers no
systematic, like-for-like comparison of these three paradigms on a single
shared evaluation pipeline. Each paper introducing a new architecture
typically compares against a small number of baselines under different
training protocols, augmentations, learning-rate schedules, and
evaluation metrics. As a result it is difficult for clinicians or
practitioners to answer a basic question: \emph{which architectural
family should I use for CAMUS-style apical 2D echocardiographic
segmentation, and what are the trade-offs?}

This paper provides that comparison on the CAMUS dataset
\citep{leclerc2019}, the standard public 2D echocardiographic
segmentation benchmark, with 500 patients across two views and two
cardiac phases. We train nine representative architectures spanning the
three paradigms above using one identical protocol --- same optimiser,
loss, schedule, augmentation, mixed-precision settings, and early-stopping
criterion --- and evaluate them on a single shared pipeline reporting
geometric, boundary, clinical, and efficiency metrics simultaneously.

\paragraph{Contributions.}
\begin{itemize}[leftmargin=*,nosep]
\item We present the first unified, like-for-like benchmark of deep
      segmentation architectures on CAMUS across the three major paradigms:
      nine architectures spanning pure CNN, hybrid CNN--Transformer, and
      pure Transformer, trained and evaluated under a single shared protocol.

\item We report \emph{multi-dimensional} evaluation, going beyond Dice
      to include HD95 in millimetres (with per-sample NIfTI pixel
      spacing), ASSD, biplane Simpson's EF MAE and Pearson $r$, ED/ES
      stratified metrics, robustness across CAMUS image-quality grades,
      and efficiency benchmarks (parameters, FLOPs, latency, peak
      memory). The full results matrix and per-patient outputs are
      released.

\item We apply pairwise Wilcoxon signed-rank tests with Bonferroni
      correction to identify which architectural differences are
      statistically significant, and produce a critical-difference
      diagram showing the top-tier band of equivalent models.

\item \textbf{We document a boundary-metric reporting error that makes
      published HD95 values on CAMUS mutually incomparable.} Computing
      the distance transform on the network's resized $256 \times 256$
      grid while attaching the native NIfTI spacing under-reports every
      distance by the resampling factor, which on CAMUS is
      $\approx 2.2\times$. The effect is large enough to invert a
      conclusion: under that convention our models score
      $1.8$--$2.9$~mm and appear to halve the published baseline, while
      the identical predictions scored at native resolution give
      $4.1$--$5.7$~mm and merely match it. We report both, and recommend
      that boundary distances be computed at acquisition resolution ---
      a concrete instance of the validation concerns raised by
      \citet{reinke2021limitations} and \citet{maierhein2024metrics}.

\item Our best architecture, TransUNet, reaches Dice $= 0.9121$,
      HD95 $= 4.08$~mm; nnU-Net reaches Dice $= 0.9099$, HD95 $= 4.27$~mm
      with $6.5\times$ fewer parameters and roughly $3\times$ faster
      inference. Boundary accuracy under our unified protocol reproduces
      the originally published nnU-Net baseline
      ($\approx 4.3$~mm) \citep{leclerc2019} rather than improving on it,
      and EF is floor-limited in the same way: the ground-truth oracle
      scores $5.44\%$ MAE under an identical pipeline against the best
      model's $6.62\%$, leaving roughly $1.2$ points attributable to
      segmentation at all.

\item We release a fully reproducible benchmark framework: training,
      evaluation, statistical analysis, and figure generation code.
\end{itemize}

\paragraph{Outline.}
Section~\ref{sec:related} reviews segmentation architectures and the
CAMUS benchmark. Section~\ref{sec:methods} describes the nine
architectures, the unified training protocol, and the evaluation
pipeline. Section~\ref{sec:results} reports the experimental results.
Section~\ref{sec:discussion} discusses architectural design choices,
trade-offs, and limitations. Section~\ref{sec:conclusion} concludes.
